Начнем с определений разных сходимостей функций распределений на прямой.

\mkdef{Сходимость в основном}{conv_main}{
    Последовательность ф.р $\{F_n, n \in \mathbb{N}\}$ на $\mathbb{R}$ называется сходящейся в основном к ф.р. $F(x)$, если $F_n \to F(x)$ для всех $x \in \mathcal{C}(F)$, где $\mathcal{C}(F)$ --- множество точек непрерывности ф.р. $F(x)$. Обозначение: $F_n \Rightarrow F$.
}

Сделаем очевидное наблюдение, что сходимость случайных величин по распределению эквивалентна сходимости в основном их функций распределения: если $\{\xi_n, n \in \mathbb{N}\}$ --- последовательность с.в., то
$$ \xi_n \xrightarrow{d} \xi \Leftrightarrow F_{\xi_n} \Rightarrow F_\xi. $$

\mkdef{Слабая сходимость}{conv_weak}{
    Пусть $\{F_n, n \in \mathbb{N}\}$ --- последовательность ф.р. на $\mathbb{R}$. Она называется слабо сходящейся к функции распределения $F(x)$, если для любой ограниченной непрерывной функции $f(x)$ выполнено
    $$ \int\limits_{\mathbb{R}} f(x) dF_n(x) \to \int\limits_{\mathbb{R}} f(x) dF(x), \quad n \to \infty. $$
    Обозначение: $F_n \xrightarrow{w} F$.
}

Для определения сходимостей вероятностных мер на метрических пространствах нам понадобится понятие борелевской $\sigma$-алгебры в подобных пространствах. Пусть $(S, \rho)$ --- метрическое пространство.

\mkdef{Борелевская сигма-алгебра}{borel_algebra}{
    Борелевской сигма-алгеброй, $\mathcal{B}(S)$, на $(S, \rho)$ называется минимальная $\sigma$-алгебра, содержащая все открытые множества из $S$.
}

При наличии сепарабельности несложно проверить следующее эквивалентное определение.

\mkprop{}{prop_sep_borel}{
    Если пространство $(S, \rho)$ сепарабельно, то $\mathcal{B}(S)$ является минимальной $\sigma$-алгеброй, содержащей все открытые шары.
}

Введем понятия слабой сходимости и сходимости в основном вероятностных мер на метрических пространствах.

\mkdef{Слабая сходимость вероятностных мер}{conv_weak_metric}{
    Пусть задано метрическое пространство $(S, \rho)$ и последовательность $\{\P_n, n \in \mathbb{N}\}$ вероятностных мер на $(S, \mathcal{B}(S))$. Будем говорить, что $\P_n$ слабо сходятся к вероятностной мере $\P$ на $(S, \mathcal{B}(S))$, если для любой ограниченной непрерывной функции $f \colon S \to \mathbb{R}$ выполнено
    $$ \lim_{n \to \infty} \int\limits_S f(x) \P_n(dx) = \int\limits_S f(x) \P(dx). $$
    Обозначение: $\P_n \xrightarrow{w} \P$.
}

\mkdef{Сходимость в основном на метрических пространствах}{conv_main_metric}{
    Последовательность вероятностных мер $\{\P_n, n \in \mathbb{N}\}$ на $(S, \rho)$ называется сходящейся в основном к мере $\P$, если
    $$ \P_n(A) \to \P(A) \quad \text{при } n \to \infty $$
    для любого $A \in \mathcal{B}(S)$ такого, что $\P(\partial A) = 0$, где через $\partial A$ обозначена граница множества $A$, т.е.
    $$ \partial A = [A] \cap [\overline{A}], $$
    $[A]$ --- это замыкание множества $A$. Обозначение: $\P_n \Rightarrow \P$.
}

Ключевую роль в теории слабой сходимости играет следующая теорема А.Д. Александрова, которая показывает, что приведенные выше виды сходимостей вероятностных мер эквивалентны, а также приводит еще несколько эквивалентных свойств.

\mkthm{Теорема Александрова}{alexandrov}{
    Пусть $\{\P_n, n \in \mathbb{N}\}$ и $\P$ --- вероятностные меры на метрическом пространстве $(S, \rho)$. Тогда следующие утверждения эквивалентны:
    \begin{enumerate}
        \item $\P_n \xrightarrow{w} \P$,
        \item $\varlimsup\limits_{n \to \infty} \P_n(F) \le \P(F)$ для любого замкнутого множества $F \subset S$,
        \item $\varliminf\limits_{n \to \infty} \P_n(G) \ge \P(G)$ для любого открытого множества $G \subset S$,
        \item $\P_n \Rightarrow \P$.
    \end{enumerate}
}

\mkproof{
    (1)$\Rightarrow$(2). Пусть $F \subset S$ --- замкнутое множество. Для произвольного $\varepsilon > 0$ введём функцию
    $$ f_\varepsilon(x) = \left( 1 - \frac{\rho(x, F)}{\varepsilon} \right)^+, \quad \text{где } \rho(x, F) = \inf_{y \in F} \rho(x, y). $$
    
    Заметим, что она неотрицательна, непрерывна и ограничена, $f_\varepsilon(x) = 1$ для всех $x \in F$ и $f_\varepsilon(x) \le \I\{x \in F^\varepsilon\}$, где $F^\varepsilon = \{x \in S \colon \rho(x, F) \le \varepsilon\}$. Тогда верна следующая цепочка соотношений:
    
    \begin{align*}
        \varlimsup_{n \to \infty} \P_n(F) &= \varlimsup_{n \to \infty} \int\limits_S \I\{x \in F\} \P_n(dx) \le \varlimsup_{n \to \infty} \int\limits_S f_\varepsilon(x) \P_n(dx) = \\
        &= \int\limits_S f_\varepsilon(x) \P(dx) \le \int\limits_S \I\{x \in F^\varepsilon\} \P(dx) = \P(F^\varepsilon),
    \end{align*}
    
    Осталось заметить, что $F^\varepsilon$ монотонно сжимается к $F$ при $\varepsilon \to 0$ в силу замкнутости $F$. В силу непрерывности вероятностной меры $\P(F^\varepsilon) \to \P(F)$. При устремлении $\varepsilon$ к нулю получаем искомое неравенство.
    \vspace{1em}
    
    (2)$\Leftrightarrow$(3) Эквивалентность второго и третьего пунктов очевидна, так как от одного к другому можно переходить, рассматривая дополнения множеств.
    \vspace{1em}
    
    (2),(3)$\Rightarrow$(4) Возьмём произвольное борелевское множество $B \in \mathcal{B}(S)$ с условием $\P(\partial B) = 0$. Введём два множества:
    
    --- $F = [B] = B \cup \partial B$ (замыкание $B$);
    
    --- $G = B \setminus \partial B$ (внутренность $B$).
    
    Заметим, что $F$ замкнуто, а $G$ открыто, причём $\P(F) = \P(B) = \P(G)$ в силу того, что $\P(\partial B) = 0$. Тогда несложно показать, что искомый предел существует и равен $\P(B)$:
    
    \begin{align*}
        \varlimsup_{n \to \infty} \P_n(B) &\le \varlimsup_{n \to \infty} \P_n(F) \le \P(F) = \P(B), \\
        \varliminf_{n \to \infty} \P_n(B) &\ge \varliminf_{n \to \infty} \P_n(G) \ge \P(G) = \P(B).
    \end{align*}
    
    Следовательно, существует $\lim\limits_{n \to \infty} \P_n(B) = \P(B)$.
    \vspace{1em}
    
    (4)$\Rightarrow$(1). Возьмём произвольную ограниченную непрерывную функцию $f \colon S \to \mathbb{R}$. Пусть $|f(x)| < M$ для всех $x \in S$. Рассмотрим множество:
    $$ D = \{t \in [-M, M] \colon \P(\{x \in S \colon f(x) = t\}) > 0\}. $$
    
    Данное множество не более, чем счётно. Теперь зафиксируем произвольное натуральное $k$ и возьмём разбиение $-M = t_0 < t_1 < \dots < t_k = M$ отрезка $[-M, M]$ такое, что ни одно из $t_i$ не содержится в $D$.
    
    Далее, построим набор множеств $\{B_i\}_{i=1}^k$ по следующему правилу:
    $$ B_i = \{x \in S \colon t_{i-1} \le f(x) < t_i\}. $$
    
    Заметим, что в силу непрерывности $f$ множество $f^{-1}((t_{i-1}, t_i))$ открыто, а $f^{-1}([t_{i-1}, t_i])$ замкнуто, поэтому $\partial B_i \subseteq f^{-1}(\{t_{i-1}\}) \cup f^{-1}(\{t_i\})$. Но оба прообраза имеют нулевую меру, поэтому граница $B_i$ тоже имеет нулевую меру. Следовательно, $\P_n(B_i) \to \P(B_i)$ при $n \to \infty$ для всех $i = 1, \dots, k$. Теперь рассмотрим следующий верхний предел:
    $$ \Delta = \varlimsup_{n \to \infty} \left| \int\limits_S f(x) \P_n(dx) - \int\limits_S f(x) \P(dx) \right|. $$
    
    Ограничим его сверху суммой трёх верхних пределов:
    \begin{align*}
        \Delta &\le \varlimsup_{n \to \infty} \left| \int\limits_S f(x) \P_n(dx) - \sum_{i=1}^k t_{i-1} \P_n(B_i) \right| + \\
        &\quad + \varlimsup_{n \to \infty} \left| \sum_{i=1}^k t_{i-1} \P_n(B_i) - \sum_{i=1}^k t_{i-1} \P(B_i) \right| + \\
        &\quad + \varlimsup_{n \to \infty} \left| \int\limits_S f(x) \P(dx) - \sum_{i=1}^k t_{i-1} \P(B_i) \right|
    \end{align*}
    
    Второй предел равен нулю. Теперь покажем, что первый (а так же и третий) предел не превосходит $\max\limits_{i=1, \dots, k} |t_i - t_{i-1}|$. Для этого заметим, что
    $$ \int\limits_S f(x) \P_n(dx) = \sum_{i=1}^k \int\limits_{B_i} f(x) \P_n(dx), $$
    и
    $$ t_{i-1} \P_n(B_i) \le \int\limits_{B_i} f(x) \P_n(dx) \le t_i \P_n(B_i). $$
    
    Следовательно,
    $$ \left| \int\limits_S f(x) \P_n(dx) - \sum_{i=1}^k t_{i-1} \P_n(B_i) \right| \le \left| \sum_{i=1}^k (t_i - t_{i-1}) \P_n(B_i) \right| \le \max_{i=1, \dots, k} |t_i - t_{i-1}|. $$
    
    Отсюда получаем, что
    $$ \Delta \le 2 \max_{i=1, \dots, k} |t_i - t_{i-1}|. $$
    
    Но этот максимум стремится к нулю при устремлении диаметра разбиения к нулю. Следовательно, $\Delta = 0$ и имеет место слабая сходимость.
}
